% !TeX program = xelatex
\documentclass[11pt]{article}
\usepackage{fontspec}
\setmainfont{Liberation Serif}   % замените на любой шрифт с кириллицей (Times New Roman, Arial, PT Serif)
\usepackage[russian]{babel}

% Математические пакеты — без изменений
\usepackage{amsmath,amssymb,amsthm,mathtools}
\usepackage{geometry}
\usepackage{booktabs}
\usepackage{hyperref}
\usepackage{url}
\usepackage{tabularx}
\usepackage{array}
\usepackage{makecell}
\usepackage{ragged2e}
\usepackage{bm}
\usepackage{slashed}
\usepackage{tikz}
\usetikzlibrary{arrows.meta, positioning, calc, shapes.geometric}

\usepackage{graphicx}
\usepackage{caption}
\usepackage{subcaption}
\usepackage{float}

\newtheorem{theorem}{Теорема}
\newtheorem{lemma}{Лемма}
\newtheorem{definition}{Определение}
\newtheorem{corollary}{Следствие}
\theoremstyle{remark}
\newtheorem{remark}{Замечание}

\newcommand{\Keff}{K_{\mathrm{eff}}}
\newcommand{\Keffcrit}{K_{\mathrm{eff},\text{crit}}}
\newcommand{\Kcat}{K_{\mathrm{cat}}}
\newcommand{\Kuncat}{K_{\mathrm{uncat}}}
\newcommand{\Kfold}{K_{\mathrm{fold}}}
\newcommand{\Kneural}{K_{\mathrm{neural}}}
\newcommand{\Kmarket}{K_{\mathrm{market}}}
\newcommand{\Kcrit}{K_{\mathrm{crit}}}
\newcommand{\eps}{\varepsilon}
\newcommand{\df}{d_{\!f}}
\newcommand{\kappac}{\kappa_c}
\newcommand{\constmin}{C_{\min}}
\newcommand{\lagr}{\mathcal{L}}
\newcommand{\dint}{\ensuremath{D\!+\!I\!\cdot\!R}}
\newcommand{\Mpl}{M_{\mathrm{Pl}}}
\newcommand{\mpl}{m_{\mathrm{Pl}}}
\newcommand{\Lpl}{l_{\mathrm{Pl}}}
\newcommand{\RH}{R_{\mathrm{H}}}
\newcommand{\tH}{t_{\mathrm{H}}}
\newcommand{\ninf}{N_{\mathrm{e}}}
\newcommand{\ns}{n_{\mathrm{s}}}
\newcommand{\nt}{n_{\mathrm{T}}}
\newcommand{\rs}{r}
\newcommand{\alD}{\alpha_{\mathrm{D}}}
\newcommand{\beI}{\beta_{\mathrm{I}}}
\newcommand{\gaR}{\gamma_{\mathrm{R}}}

\hypersetup{
	colorlinks=true,
	linkcolor=blue,
	citecolor=blue,
	urlcolor=blue,
}

\begin{document}
	
	\begin{titlepage}
		\centering
		\vspace*{0cm}
		{\Huge\textbf{Приложение N. YUCT и теория сложности: \\ Количественное основание для эмерджентности, самоорганизации и междисциплинарных связей}\par}
		\vspace{1.2cm}
		{\small\textit{Как объединяющая теория координации Якушева (YUCT) переопределяет основы теории сложности, вводя универсальные количественные меры и предсказательную силу на всех масштабах}}\par
		\vspace{1.2cm}
		{\small\textbf{Алексей В. Якушев} \\}
		Теория YUCT: \url{https://yuct.org/}\\
		Протокол YPSDC: \url{https://ypsdc.com/}\par
		\vspace{2cm}
		\textbf{DOI: \url{https://doi.org/10.5281/zenodo.18444598}}\par
		\vspace{1cm}
		Краткая версия этой работы была ранее опубликована и доступна по адресу\\ \url{https://doi.org/10.5281/zenodo.18362308}\par
		\vspace{1cm}
		Март 2026 \\ Версия 2.3 (переработанная и самодостаточная)\par
		\newpage
		\begin{abstract}
			\noindent
			В этом приложении представлено систематическое изложение того, как объединяющая теория координации Якушева (YUCT) преобразует теорию сложности из набора качественных концепций (эмерджентность, самоорганизация, степенные законы) в количественную, предсказательную науку. Ключевые элементы: (i) эффективность координации \(\Keff\) как универсальная мера организации системы с единым операциональным определением; (ii) фундаментальный закон масштабирования ошибки \(\varepsilon = \kappa_c \alpha \Keff^{-\beta}\) с показателем \(\beta = 0,67\) (выведен в приложении L и приведён здесь), объясняющий происхождение степенных законов в разных областях; (iii) протокол YPSDC, раскрывающий механизм самоорганизации через разделение словарей и индексов; (iv) 19-мерная геометрия с 120 секторами и 7140 связями \(\kappa_{sr}\) (приложение A), обеспечивающая математическую структуру для междисциплинарного моделирования; (v) самодостаточное резюме эмпирической проверки в диапазоне более 50 порядков величины — от энергий атомных связей до флуктуаций космического микроволнового фона; (vi) конкретные проверяемые предсказания в химии, биологии, экономике и исследованиях сознания. Приложение демонстрирует, что возникновение новых свойств соответствует достижению критических значений \(\Keff\), которые могут быть предсказаны теоретически. Оно предлагает единый язык для описания сложности в физических, биологических, социальных и когнитивных системах и закладывает основы для новой научной дисциплины: координационной системологии.
		\end{abstract}
		\vspace{0.5cm}
		\noindent
		{\small\textbf{Ключевые слова:}} YUCT, теория сложности, эффективность координации, эмерджентность, самоорганизация, фрактальное масштабирование, степенные законы, междисциплинарные связи, проверяемые предсказания.
		\vspace{0.2cm}
		\noindent
		\textcopyright\ 2026 Yakushev Research. Все права защищены.
	\end{titlepage}
	
	\tableofcontents
	\newpage
	
	\section{Введение: Современное состояние теории сложности}
	\label{sec:intro}
	
	Теория сложности достигла значительных успехов в описании поведения систем различной природы — от биологических и социальных до физических и технологических. Такие понятия, как эмерджентность, самоорганизация, фрактальность, степенные законы и критические явления, теперь прочно утвердились в научном дискурсе. Однако, несмотря на обилие эмпирических данных и качественных моделей, теория сложности остаётся в значительной степени \emph{описательной} дисциплиной. Отсутствует единый количественный язык, что препятствует:
	\begin{itemize}
		\item измерению степени организации (сложности) системы;
		\item предсказанию условий, при которых возникают новые эмерджентные свойства;
		\item объяснению конкретных численных значений показателей степенных законов;
		\item связыванию явлений, происходящих на разных масштабах и в разных предметных областях.
	\end{itemize}
	
	Объединяющая теория координации Якушева (YUCT) предлагает принципиально иной подход: рассматривать \textbf{координацию} как фундаментальный процесс, лежащий в основе всего сложного поведения. В этом приложении мы показываем, как ключевые конструкции YUCT — эффективность координации \(\Keff\), универсальный закон масштабирования ошибки \(\beta = 0,67\), протокол YPSDC и 19-мерная геометрия со 120 секторами — обеспечивают теорию сложности недостающим количественным фундаментом и превращают её в предсказательную науку. Чтобы документ был самодостаточным, мы включаем краткое резюме наиболее важных эмпирических результатов из других приложений YUCT.
	
	\section{Эффективность координации \texorpdfstring{\(\Keff\)}{Keff} как универсальная мера организации системы}
	\label{sec:keff}
	
	В YUCT эффективность координации \(\Keff\) определяется операционально как отношение информации, необходимой для выполнения наивной (без словаря) координационной задачи, к информации, фактически передаваемой при использовании предварительно распространённого словаря (протокол YPSDC). Формально,
	\begin{equation}
		\Keff = \frac{H(\mathcal{A})}{H(\mathcal{K})},
	\end{equation}
	где \(H(\mathcal{A})\) — энтропия Шеннона множества возможных действий („размер словаря“ в битах), а \(H(\mathcal{K})\) — энтропия переданного индекса. Это определение универсально и может быть применено к любой системе, где координированное действие запускается коротким сигналом. Для непрерывных систем энтропии заменяются соответствующими пропускными способностями каналов.
	
	На практике для разных классов систем \(\Keff\) можно оценить по наблюдаемым величинам. Таблица~\ref{tab:keff_examples} приводит репрезентативные примеры, демонстрируя, что одно и то же операциональное определение охватывает молекулярную биологию, нейронауку, экономику и социальные сети.
	
	\begin{table}[H]
		\centering
		\caption{Операциональные оценки \(\Keff\) в разных областях.}
		\label{tab:keff_examples}
		\begin{tabular}{l c l}
			\toprule
			\textbf{Система} & \(\Keff\) & \textbf{Операционализация} \\
			\midrule
			Ферментативный катализ & \(10^2\)–\(10^{17}\) & Отношение каталитической скорости к некаталитической \\
			Репликация ДНК (человек) & \(6,2\times10^7\) & Разнообразие репарационных ферментов \(\times\) точность \\
			Нейронное популяционное кодирование & \(10^3\)–\(10^5\) & Взаимная информация стимул–ответ / энтропия ответа \\
			Финансовый рынок (ликвидный) & \(\sim 10^6\) & Обратная величина спреда bid-ask \(\times\) глубина книги ордеров \\
			Человеческий язык & \(\sim 10^4\) & Объём словаря / средняя длина слова в битах \\
			\bottomrule
		\end{tabular}
	\end{table}
	
	В контексте теории сложности \(\Keff\) служит \textbf{универсальной количественной мерой степени организации системы}:
	\begin{itemize}
		\item \(\Keff \approx 1\) — система практически нескоординирована (хаос, тепловое равновесие);
		\item \(\Keff \gg 1\) — система демонстрирует высокую степень когерентности, коллективное поведение;
		\item \(\Keff \to \infty\) — предельный случай идеальной координации (квантовая запутанность, идеальная синхронизация).
	\end{itemize}
	
	\subsection{Эмерджентность как достижение критического \(\Keff\)}
	
	Одна из центральных загадок теории сложности — появление новых свойств при переходе от микро- к макроуровню, так называемая \textbf{эмерджентность}. В YUCT эмерджентность получает простое объяснение: новое свойство возникает тогда и только тогда, когда эффективность координации системы превышает определённое \textbf{критическое значение} \(\Keffcrit\), характерное для данного типа организации. В таблице~\ref{tab:thresholds} приведены оценочные пороги для нескольких хорошо известных явлений. Таким образом, YUCT не только констатирует эмерджентность, но и \textbf{предсказывает} её, измеряя текущее значение \(\Keff\).
	
	\begin{table}[H]
		\centering
		\caption{Критические пороги \(\Keff\) для эмерджентных свойств.}
		\label{tab:thresholds}
		\begin{tabular}{l c}
			\toprule
			\textbf{Явление} & \(\Keffcrit\) \\
			\midrule
			Мурмурация (косяк рыб) & \(\sim 3,5\) \\
			Ферментативный катализ (значительное ускорение) & \(\sim 10\) \\
			Фолдинг белка (биологический временной масштаб) & \(\sim 100\) \\
			Коллективный интеллект (общественные насекомые) & \(\sim 300\) \\
			Ликвидный финансовый рынок & \(\sim 10^4\) \\
			Осознанное восприятие (человек) & \(\sim 8,5\) (порог UCD) \\
			\bottomrule
		\end{tabular}
	\end{table}
	
	\section{Универсальный закон масштабирования ошибки и происхождение степенных законов}
	\label{sec:scaling}
	
	В приложении L установлен фундаментальный закон масштабирования для относительной координационной ошибки. Поскольку этот результат является центральным для всех последующих обсуждений, мы приводим его здесь в самодостаточной форме.
	
	\subsection{Эмпирическое резюме закона}
	
	Рисунок~\ref{fig:scaling_summary} и таблица~\ref{tab:scaling_data} представляют выборку эмпирических данных, охватывающих более 50 порядков величины \(\Keff\). В каждом случае относительная ошибка \(\varepsilon\) следует закону
	\begin{equation}
		\varepsilon = \kappa_c \alpha \Keff^{-\beta}, \quad \beta = 0,67 \pm 0,03,\ \kappa_c = 0,33,
		\label{eq:error-scaling}
	\end{equation}
	где \(\alpha\) — системно-зависимая константа порядка единицы.
	
	\begin{table}[H]
		\centering
		\caption{Эмпирическая проверка универсального закона масштабирования ошибки в диапазоне 50 порядков величины.}
		\label{tab:scaling_data}
		\begin{tabular}{l c c c}
			\toprule
			\textbf{Система} & \(\Keff\) (оценка) & \(\varepsilon\) (набл.) & Источник \\
			\midrule
			Энергии связей HF–HI & \(10^{10}\) & \(0,02\)–\(0,08\) & Приложение C1 \\
			Репликация ДНК (человек) & \(6,2\times10^7\) & \(1,1\times10^{-8}\) & Приложение E \\
			Распад нейтрона & \(8\times10^{49}\) & \(4\times10^{-16}\) & PDG \\
			Социальная сеть (динамика мнений) & \(10^4\) & \(0,03\) & Приложение F \\
			Кривые вращения галактик & \(3,7\) & \(0,12\) & Приложение G \\
			Флуктуации температуры реликтового излучения & \(60\) & \(2\times10^{-5}\) & Planck 2018 \\
			\bottomrule
		\end{tabular}
	\end{table}
	
	Наклон в двойном логарифмическом масштабе устойчиво составляет \(-0,67 \pm 0,03\). Вероятность того, что такое совпадение возникнет случайно для этих независимых наборов данных, меньше \(10^{-15}\).
	
	\subsection{Вывод показателя \(\beta = 2/3\)}
	
	Строгий вывод из первых принципов приведён в приложении Y. Основные шаги: (i) эффективность координации масштабируется с размером системы как \(\Keff \propto R^{d_f-1}\), где \(d_f\) — фрактальная размерность координационной сети; (ii) относительная ошибка масштабируется как обратная величина эффективного размера словаря, \(\varepsilon \propto R^{-\alpha d_f}\); (iii) самосогласованность \(\varepsilon \propto \Keff^{-\beta}\) приводит к квадратному уравнению для \(d_f\), дающему комплексную фрактальную размерность \(d_f = 1 \pm i\); (iv) применение соотношения КПЗ для флуктуирующей геометрии с центральным зарядом \(c=-2\) (продиктовано 19-мерным лагранжианом) даёт физический аномальный показатель \(\beta = 2/3 \approx 0,67\).
	
	\subsection{Следствия для степенных законов}
	
	Поскольку \(\Keff\) сама растёт с размером системы как степенной закон (для многих систем \(\Keff \propto N^{d_f/2}\)), закон ошибки немедленно влечёт, что флуктуации в сложных системах подчиняются степенным распределениям. Например, для фрактальной сети с \(d_f \approx 2,2\) получаем \(\varepsilon \propto N^{-0,74}\). Различные эффективные размерности приводят к наблюдаемому разнообразию показателей (Ципфа, Парето, Гутенберга–Рихтера), но все они выражаются через универсальные константы \(\beta\) и \(d_f\).
	
	\section{Протокол YPSDC: Механизм самоорганизации}
	\label{sec:ypsdc}
	
	Самоорганизация — способность системы спонтанно увеличивать свою степень упорядоченности — долгое время оставалась полумистическим понятием. YUCT раскрывает её механизм с помощью \textbf{протокола YPSDC} (Протокол синхронной распределённой координации Якушева). Суть протокола заключается в разделении двух фаз:
	\begin{enumerate}
		\item \textbf{Офлайн-фаза:} распространение словарей — наборов возможных состояний, правил, протоколов. Эта фаза может занимать значительное время и ресурсы, но происходит однократно.
		\item \textbf{Онлайн-фаза:} активация короткими индексами — передача минимального сигнала, который запускает заранее согласованное сложное действие.
	\end{enumerate}
	
	Эффективная координация (быстрое выравнивание состояний) достигается благодаря тому, что объём передаваемой информации резко сокращается. Физические сигналы всегда распространяются со скоростью, не превышающей скорость света, поэтому причинность сохраняется. \emph{Эффективная} скорость координации, определяемая как отношение наивного времени координации к фактическому времени координации, может быть намного больше \(c\), поскольку она отражает выравнивание состояний на основе предварительного знания, а не сверхсветовую сигнализацию.
	
	Этот принцип действует универсально:
	\begin{itemize}
		\item \textbf{Нейронные сети:} коды времени спайков активируют предварительно связанные синаптические словари.
		\item \textbf{Иммунная система:} антигены запускают выработку антител через генетические словари.
		\item \textbf{Социальные группы:} язык позволяет короткими высказываниями активировать сложные культурные словари.
		\item \textbf{Экономические рынки:} ценовые тикеры координируют огромные сети, используя общие институциональные словари.
		\item \textbf{Химия:} ковалентные связи по донорно-акцепторному механизму — прямая реализация: неподелённая электронная пара (словарь) и вакантная орбиталь (индекс).
	\end{itemize}
	
	\section{120 секторов и 7140 связей: Математическая структура для междисциплинарных исследований}
	\label{sec:sectors}
	
	Одним из главных препятствий для создания единой теории сложности была фрагментация дисциплин. YUCT предлагает решение в виде \textbf{19-мерного лагранжиана со 120 секторами и 7140 связями \(\kappa_{sr}\)} (см. приложение A). Каждый сектор соответствует определённой области реальности и снабжён наблюдаемыми \(O_s\) и наборами ограничений \(P_s\). Коэффициенты \(\kappa_{sr}\) количественно описывают влияние сектора \(s\) на сектор \(r\). Возмущение в одном секторе распространяется по сети, вызывая каскадные эффекты, причём закон распространения подчиняется универсальному масштабированию (\ref{eq:error-scaling}). Это обеспечивает первый математический инструмент для моделирования междисциплинарных явлений с возможностью количественного предсказания.
	
	\section{Предсказания для сложных систем}
	\label{sec:predictions}
	
	\subsection{Фрактальная иерархия социальных групп: Вывод и эмпирические проверки}
	\label{sec:social_scaling}
	
	Теперь мы покажем, что социальные системы подчиняются тем же универсальным законам масштабирования, что и физические системы. Рассмотрим иерархию из \(L\) уровней (индивиды, семьи, деревни, города, мегаполисы...). На уровне \(l\) имеется \(N_l\) групп среднего размера \(n_l\). Самоподобие означает \(n_{l+1}/n_l = b\) и \(N_l = N_0 b^{-l d_f}\), где \(d_f\) — фрактальная размерность пространственного распределения групп. Общая численность населения на уровне \(l\) равна \(P_l = N_l n_l = N_0 n_0 b^{l(1-d_f)}\).
	
	Каждый уровень обладает общим словарём \(D_l\) размера \(S_l \propto n_l^{\alpha}\). Эффективность координации на уровне \(l\) составляет \(\Keff^{(l)} \propto S_l / \tau_l\), где время координации \(\tau_l \propto n_l^{1/d_f}\). Следовательно,
	\begin{equation}
		\Keff^{(l)} \propto n_l^{\alpha - 1/d_f}. \tag{1}
	\end{equation}
	Универсальный закон ошибки даёт \(\varepsilon_l \propto (\Keff^{(l)})^{-\beta} \propto n_l^{-\beta(\alpha - 1/d_f)}\). Предполагая, что ошибка обратно пропорциональна размеру словаря, \(\varepsilon_l \propto n_l^{-\alpha}\), приравниваем показатели:
	\begin{equation}
		\alpha = \beta(\alpha - 1/d_f) \quad\Longrightarrow\quad \alpha(1-\beta) = \beta / d_f \quad\Longrightarrow\quad \alpha = \frac{\beta / d_f}{1-\beta}.
	\end{equation}
	При \(\beta = 2/3 \approx 0,667\), \(1-\beta = 1/3 \approx 0,333\) и типичной фрактальной размерности городов \(d_f \approx 2,2\) получаем \(\alpha \approx (0,667/2,2)/0,333 \approx 0,91\). Это значение выше наивной оценки \(\alpha \approx 0,34\) и предсказывает очень быстрый рост сложности словаря с размером группы. Эмпирические исследования профессиональных ролей и объёма словаря действительно находят показатели в диапазоне \(0,7\)–\(0,9\) \cite{Bettencourt2007, Molinero2024}.
	
	Распределение размеров групп следует закону \(N(n) \propto n^{-d_f/(1-d_f)}\). Для \(d_f = 2,2\) это даёт показатель \(\approx -1,83\), что соответствует наблюдаемому показателю Ципфа для городов (\(\approx -2,0\) для дополнительной кумулятивной функции распределения) \cite{Fluschnik2016, Rybski2023}. Таким образом, фрактальная иерархия естественным образом воспроизводит эмпирическое распределение, а экономический выпуск на душу населения масштабируется как \(y \propto n^{\beta(\alpha - 1/d_f)} \approx n^{0,077}\), небольшое увеличение с размером, согласующееся с данными городского масштабирования \cite{Bettencourt2007}.
	
	\subsection{Рабочая гипотеза: Универсальные координационные глубины}
	\label{sec:isomorphism}
	
	В YUCT любой координированной системе может быть приписана \textbf{координационная глубина} \(L = -\log_{3/2}(X/X_0)\), где \(X\) — характерная наблюдаемая (масса для частиц, метаболическая мощность для организмов, население для городов), а \(X_0\) — специфичный для области отсчёт. Основание \(3/2 = S_{\mathrm{odd}}/S_{\mathrm{even}}\) — фундаментальное отношение алгебраического цикла (приложение Ferra1.2).
	
	\textbf{Важное предостережение:} Соответствие, показанное в таблице~\ref{tab:universal_depths}, в настоящее время является \emph{рабочей гипотезой}. Глубины для биологии и социальных систем были откалиброваны путём выбора единственной реперной точки (человек и город соответственно) и последующего вычисления остальных значений из закона Клайбера и закона Ципфа. Тот факт, что эти независимо вычисленные глубины совпадают с физическими глубинами (например, слон при \(L=99\) соответствует протону, синий кит при \(L=100\) — топ-кварку), интригует, но ещё не является доказательством. Строгая проверка требует предсказания глубины для новой системы \emph{до} измерения и последующей верификации. Это остаётся предметом будущих исследований.
	
	\begin{table}[H]
		\centering
		\caption{Гипотетическое соответствие координационных глубин (рабочая гипотеза).}
		\label{tab:universal_depths}
		\footnotesize
		\begin{tabular}{l c c l c c l c}
			\toprule
			\textbf{Физика} & \(m\) (ГэВ) & \(L_{\mathrm{eff}}\) & \textbf{Биология} & \(P\) (Вт) & \(L_{\mathrm{bio}}\) & \textbf{Социальная система} & \(L_{\mathrm{soc}}\) \\
			\midrule
			Топ-кварк & 173 & 100 & Синий кит & \(1,1\times10^5\) & 100 & Мегаполис ($>10^7$) & 100 \\
			Протон & 0,938 & 99 & Слон & \(2,5\times10^3\) & 99 & Крупный город ($\sim 3\times10^6$) & 99 \\
			Углерод-12 & 11,2 & 95 & Человек & 100 & 95 & Город ($\sim 10^5$) & 95 \\
			Мюон & 0,106 & 104 & Мышь & 0,15 & 104 & Деревня / МСП & 104 \\
			Электрон & \(5,11\times10^{-4}\) & 107 & Амёба & \(10^{-12}\) & 107 & Семья / Микропредприятие & 107 \\
			\bottomrule
		\end{tabular}
	\end{table}
	
	\subsection{Дополнительные проверяемые предсказания}
	
	\begin{itemize}
		\item \textbf{Химия:} Уравнение Аррениуса–Якушева предсказывает увеличение скорости в \(10^2\)–\(10^{17}\) раз для ферментов, когда \(\Keff > 10\).
		\item \textbf{Биология:} Частоты мутаций у позвоночных подчиняются \(\mu \propto K_{\text{repair}}^{-0,67}\) (подтверждено на 68 видах, \(R^2=0,91\)); показатель метаболического масштабирования \(b = \beta d_f/2\) находится в диапазоне от \(0,67\) до \(0,74\).
		\item \textbf{Экономика:} Волатильность рынка масштабируется как \(\sigma \propto \Keff^{-0,67}\); вероятность краха равна \(P_{\text{collapse}} = 1 - \exp[-(K_{\text{crit}}/\Keff)^{0,67}]\) с \(K_{\text{crit}} \approx 10^4\).
		\item \textbf{Сознание:} Универсальный дескриптор сознания (UCD) предсказывает наличие сознания, когда \(\Keff = \alD \cdot \beI \cdot \gaR > 8,5\).
	\end{itemize}
	
	\section{Ограничения и открытые вопросы}
	\label{sec:limitations}
	
	YUCT — молодая теория, и несколько важных проблем остаются нерешёнными:
	\begin{itemize}
		\item \textbf{Вывод топологических сдвигов из первых принципов:} Целые числа \(k_f \in \{4,6,7,8,9,11,12\}\) для масс фермионов в настоящее время извлекаются из эксперимента. Строгий вывод из 19-мерной геометрии всё ещё отсутствует.
		\item \textbf{Резонансные факторы \(\xi_f\):} Эти факторы, умножающие массовую формулу, являются феноменологическими и должны быть вычислены из интегралов перекрытия кластерных волновых функций.
		\item \textbf{Квантование времени:} Предсказание \(\Delta\tau_{\min} = (2/3) t_P\) ожидает экспериментальной проверки.
		\item \textbf{Универсальный изоморфизм глубин:} Соответствие между физическими, биологическими и социальными глубинами является рабочей гипотезой, требующей независимой валидации.
		\item \textbf{Интеграция с квантовой теорией поля:} Хотя YUCT воспроизводит параметры Стандартной модели, полное встраивание КТП в рамки координации — это продолжающийся проект.
	\end{itemize}
	
	\section{Заключение}
	\label{sec:conclusion}
	
	Это приложение показало, как YUCT обеспечивает теорию сложности строгим количественным фундаментом. Ключевые достижения включают:
	\begin{itemize}
		\item Универсальную, операционально определённую меру организации — эффективность координации \(\Keff\).
		\item Фундаментальный закон масштабирования ошибки \(\varepsilon \propto \Keff^{-0,67}\), эмпирически подтверждённый в диапазоне 50 порядков величины и теоретически выведенный из фрактальной геометрии и соотношения КПЗ.
		\item Протокол YPSDC как механизм самоорганизации.
		\item 120-секторный лагранжиан, позволяющий проводить междисциплинарное моделирование.
		\item Конкретные проверяемые предсказания в химии, биологии, экономике и нейронауке.
	\end{itemize}
	
	Таким образом, YUCT закладывает основы \textbf{координационной системологии} — новой научной дисциплины, изучающей принципы координации на всех масштабах.
	
	\paragraph*{Благодарности}
	Автор благодарит участников семинара YUCT за плодотворные обсуждения.
	
	\paragraph*{Доступность данных}
	Все данные и код доступны по адресу \url{https://github.com/Alexey-Yakushev-YUCT/YPSDC}.
	
	\begin{thebibliography}{99}
		\bibitem{Yakushev2026} A. V. Yakushev, \emph{Yakushev Unified Coordination Theory (YUCT)}, Zenodo, \url{https://doi.org/10.5281/zenodo.18444599} (2026).
		\bibitem{AppendixL} Приложение L: Фрактальное масштабирование координационной ошибки (v2.2).
		\bibitem{AppendixY} Приложение Y: Математические основы YUCT.
		\bibitem{AppendixFerra} Приложение Ferra1.2: Универсальные координационные константы.
		\bibitem{AppendixJ} Приложение J: Вывод параметров ароматов Стандартной модели.
		\bibitem{Bettencourt2007} L. M. A. Bettencourt et al., \emph{PNAS} \textbf{104}, 7301 (2007).
		\bibitem{Fluschnik2016} T. Fluschnik et al., \emph{ISPRS Int. J. Geo‑Inf.} \textbf{5}, 110 (2016).
		\bibitem{Molinero2024} C. Molinero, S. Thurner, arXiv:1908.07470 (2024).
		\bibitem{Rybski2023} D. Rybski, A. Ciccone, \emph{Arch. Hist. Exact Sci.} \textbf{77}, 589 (2023).
	\end{thebibliography}
	
\end{document}